%% optmodel-doc.tex -- manual for the optmodel package
%% Copyright 2026 Krtiiik
%% Distributed under the LaTeX Project Public License, version 1.3c or later.
\documentclass[a4paper,11pt]{article}
\usepackage[T1]{fontenc}
\usepackage{lmodern}
\usepackage{amssymb}
\usepackage{optmodel}
\usepackage{booktabs}
\usepackage{array}

\newcommand\cs[1]{\texttt{\textbackslash #1}}
\newcommand\pkg[1]{\textsf{#1}}
\makeatletter
\def\verbatim@font{\normalfont\ttfamily\footnotesize}
\makeatother
\setlength\parindent{0pt}
\setlength\parskip{0.6\baselineskip}

\title{The \pkg{optmodel} package\\[2pt]\large typesetting optimization models}
\author{Krtiiik}
\date{Version 1.0 \quad 2026/10/06}

\begin{document}
\maketitle

\section{Introduction}

\pkg{optmodel} provides one environment, \texttt{optmodel}, that lays out a
mathematical optimization model: a keyword (\verb|\min|, \verb|s.t.|), the
left-hand sides, relations and right-hand sides of the constraints aligned in
columns, an optional quantifier column (\verb|\forall t|), optional per-line
numbers or tags, and an optional tag for the whole model.
The block is centred (or indented) as a unit and never exceeds the line width.

All public names carry the prefix \texttt{opt}, so the package is safe to load
next to other packages.
It needs a \LaTeX{} kernel from 2020-10-01 or later and loads \pkg{amsmath},
\pkg{mathtools} and \pkg{array}.

\section{Syntax}

\begin{verbatim}
\usepackage{optmodel}
...
\begin{optmodel}[<model tag>][<label key>]
  <keyword> & <lhs> & <rel> & <rhs> & <quantifier> & <line tag> \\
  ...
\end{optmodel}
\end{verbatim}

Every row has up to six columns, separated by \verb|&|:
(1)~keyword, (2)~left-hand side (right-aligned), (3)~relation, (4)~right-hand
side, (5)~quantifier, (6)~line number or tag (right-aligned at the margin).
Columns 2--5 are in display math mode; column~1 is too, so use \verb|\text{s.t.}|
for words.
Missing trailing cells may be omitted.
Rows are separated by \verb|\\|; the space between rows is \cs{optrowsep} unless
you give the usual optional argument, \verb|\\[2pt]|.

Both optional arguments of the environment are optional:

\begin{center}\small
\begin{tabular}{@{}>{\ttfamily}l >{\raggedright\arraybackslash}p{0.5\textwidth}@{}}
\toprule
\textnormal{Call} & \textnormal{Effect}\\
\midrule
\verb|\begin{optmodel}|          & no model tag\\
\verb|\begin{optmodel}[PP2]|     & model tag \textnormal{(PP2)} at the right margin\\
\verb|\begin{optmodel}[]|        & auto-numbered model tag, uses the \texttt{equation} counter\\
\verb|\begin{optmodel}[..][key]| & additionally \cs{label}\texttt{\{key\}}, so \cs{ref}\texttt{\{key\}} yields the tag\\
\bottomrule
\end{tabular}
\end{center}

\section{Modes}

All examples use the same small model unless noted.

\subsection{No label}

\begin{verbatim}
\begin{optmodel}
  \min        & \optobjective{c^{\top} x} \\
  \text{s.t.} & A x & \le & b \\
              & x   & \ge & 0
\end{optmodel}
\end{verbatim}
\begin{optmodel}
  \min        & \optobjective{c^{\top} x} \\
  \text{s.t.} & A x & \le & b \\
              & x   & \ge & 0
\end{optmodel}

\subsection{Custom label}

\begin{verbatim}
\begin{optmodel}[LP]
  ...
\end{optmodel}
\end{verbatim}
\begin{optmodel}[LP]
  \min        & \optobjective{c^{\top} x} \\
  \text{s.t.} & A x & \le & b \\
              & x   & \ge & 0
\end{optmodel}

\subsection{Automatic label and a label key}

An empty first argument takes the next equation number; the second argument is
a \cs{label} key.

\begin{verbatim}
\begin{optmodel}[][m:lp]
  ...
\end{optmodel}
Model~\ref{m:lp} is a linear program.
\end{verbatim}
\begin{optmodel}[][m:lp]
  \min        & \optobjective{c^{\top} x} \\
  \text{s.t.} & A x & \le & b \\
              & x   & \ge & 0
\end{optmodel}
Model~\ref{m:lp} is a linear program.

\subsection{Per-line numbers: \cs{optnum}}

Put \cs{optnum} in column~6 to number a line with the next equation number.
Labels work as usual.

\begin{verbatim}
\begin{optmodel}
  \min        & \optobjective{c^{\top} x} \\
  \text{s.t.} & A x & \le & b & & \optnum\label{c:cap} \\
              & x   & \ge & 0 & & \optnum
\end{optmodel}
Constraint~\ref{c:cap} is the capacity constraint.
\end{verbatim}
\begin{optmodel}
  \min        & \optobjective{c^{\top} x} \\
  \text{s.t.} & A x & \le & b & & \optnum\label{c:cap} \\
              & x   & \ge & 0 & & \optnum
\end{optmodel}
Constraint~\ref{c:cap} is the capacity constraint.

\subsection{Per-line tags: \cs{opttag}, with a model tag}

\cs{opttag}\verb|{<text>}| prints a fixed tag; \cs{label} then refers to that
text.
A row consisting of a single \cs{optwholerow} spans columns 2--4 and is left
aligned; it is meant for domains.

\begin{verbatim}
\begin{optmodel}[PP2][m:pp2]
  \min        & \optobjective{\sum_t c_t x_t} \\
  \text{s.t.} & \sum_t A^{\mathrm{I}}_t x_t & \le & b^{\mathrm{I}}
    & & \opttag{I}\label{c:I} \\
  & A^{\mathrm{II}}_t x_t + B_t y_t & \le & b^{\mathrm{II}}_t,
    & \forall t & \opttag{II}\label{c:II} \\
  & \optwholerow{x_t \in \mathbf{X}_t,\ y_t \in \mathbf{Y}_t} & &
\end{optmodel}
Model~\ref{m:pp2} has linking constraints~\ref{c:I} and~\ref{c:II}.
\end{verbatim}
\begin{optmodel}[PP2][m:pp2]
  \min        & \optobjective{\sum_t c_t x_t} \\
  \text{s.t.} & \sum_t A^{\mathrm{I}}_t x_t & \le & b^{\mathrm{I}}
    & & \opttag{I}\label{c:I} \\
  & A^{\mathrm{II}}_t x_t + B_t y_t & \le & b^{\mathrm{II}}_t,
    & \forall t & \opttag{II}\label{c:II} \\
  & \optwholerow{x_t \in \mathbf{X}_t,\ y_t \in \mathbf{Y}_t} & &
\end{optmodel}
Model~\ref{m:pp2} has linking constraints~\ref{c:I} and~\ref{c:II}.

\subsection{Centred versus indented}

By default the block is centred in the line (switch \cs{optcentertrue}).
After \cs{optcenterfalse} it is instead indented by \cs{optindent} from the left
margin.
The indent is reduced when the model would otherwise run into the tags.
Both switches and lengths are ordinary \TeX{} assignments, so keep them in a group.

\begin{verbatim}
\begingroup
  \optcenterfalse \setlength\optindent{3em}
  \begin{optmodel}[LP] ... \end{optmodel}
\endgroup
\end{verbatim}
\begingroup
  \optcenterfalse \setlength\optindent{3em}
  \begin{optmodel}[LP]
    \min        & \optobjective{c^{\top} x} \\
    \text{s.t.} & A x & \le & b \\
                & x   & \ge & 0
  \end{optmodel}
\endgroup

\subsection{Long objective}

\cs{optobjective} spans columns 2--4 but adds nothing to their widths: the
constraints alone determine the layout, and a long objective starts at the
left edge of the left-hand-side column and simply extends to the right of the
constraints.
The line numbers below belong to the constraints only.

\begin{verbatim}
\begin{optmodel}
  \min        & \optobjective{\sum_{t=1}^{T}\sum_{k=1}^{K}
    \bigl(c_{tk}x_{tk}+f_{tk}y_{tk}+h_{tk}s_{tk}\bigr)
    + \sum_{t=1}^{T} p_t u_t} \\
  \text{s.t.} & s_{t-1,k}+x_{tk}-s_{tk} & = & d_{tk},
    & \forall t,k & \optnum \\
  & x_{tk} & \le & M y_{tk}, & \forall t,k & \optnum
\end{optmodel}
\end{verbatim}
\begin{optmodel}
  \min        & \optobjective{\sum_{t=1}^{T}\sum_{k=1}^{K}
    \bigl(c_{tk}x_{tk}+f_{tk}y_{tk}+h_{tk}s_{tk}\bigr)
    + \sum_{t=1}^{T} p_t u_t} \\
  \text{s.t.} & s_{t-1,k}+x_{tk}-s_{tk} & = & d_{tk},
    & \forall t,k & \optnum \\
  & x_{tk} & \le & M y_{tk}, & \forall t,k & \optnum
\end{optmodel}

\section{Parameters}

The spacing parameters are ordinary lengths and the centring is a switch.
They may be changed at any point and apply to the models that follow in the
same group.

\begin{center}\small
\begin{tabular}{@{}>{\ttfamily\textbackslash}l l >{\raggedright\arraybackslash}p{0.4\textwidth}@{}}
\toprule
\multicolumn{1}{@{}l}{Name} & Default & Meaning\\
\midrule
optkeysep   & 1em   & gap between keyword and left-hand side\\
optquantsep & 1em   & gap between right-hand side and quantifier\\
optnumsep   & 2em   & minimum gap before the line numbers\\
opttagsep   & 1.5em & gap between line numbers and the model tag\\
optrowsep   & 4pt   & extra space between rows\\
optindent   & 2em   & left indent when \cs{optcenterfalse} is in effect\\
\bottomrule
\end{tabular}
\end{center}

\begin{center}\small
\begin{tabular}{@{}>{\ttfamily\textbackslash}l >{\raggedright\arraybackslash}p{0.6\textwidth}@{}}
\toprule
\multicolumn{1}{@{}l}{Command} & Meaning\\
\midrule
optobjective\{..\} & objective row, starts at the left-hand-side column, adds no width\\
optwholerow\{..\}  & free row across columns 2--4, left aligned, adds width\\
optnum             & auto-numbered line tag (column 6)\\
opttag\{..\}       & fixed-text line tag (column 6)\\
optcentertrue      & centre the block (default)\\
optcenterfalse     & indent the block by \cs{optindent}\\
\bottomrule
\end{tabular}
\end{center}

\section{Limitations}

\begin{itemize}
\item The body is read as an argument and typeset twice (once to measure it).
  Verbatim material (\cs{verb}, \texttt{listings}) is not allowed inside, and
  commands with side effects other than \cs{label}, \cs{optnum} and \cs{opttag}
  run twice.
\item A model is a single unbreakable block; it will not be split across pages.
\item An objective is not line-broken. One that is wider than the space to the
  right of its column runs past the margin or into the model tag, and \TeX{}
  does \emph{not} report this as an overfull box. Check long objectives by eye
  and split them by hand.
\item The layout has the fixed six columns described above.
\item With \pkg{hyperref} the package compiles cleanly and \cs{ref} prints the
  right text, but a \cs{label} after \cs{opttag} links to the enclosing
  section, not to the line. Model tags and \cs{optnum} lines link correctly.
\item The environment builds on \texttt{equation*}; do not nest it in other
  display environments.
\end{itemize}

\end{document}
